%------------------------------------------------------------------------------%
% \chapter{Plano Cartesiano e Equações}
%------------------------------------------------------------------------------%

%------------------------------------------------------------------------------%
\begin{exercicios}{Revisão}
    
    \begin{exercicio}
        Qual a equação da circunferência de raio~$3$ cujo centro é a intersecção das retas $x -2y +1 = 0$ e $x + y - 2 = 0$.
    \end{exercicio}
        
    \begin{exercicio}
        Construa a equação da reta $s$ que passa pelo ponto $(2,-2)$ e é perpendicular a reta $r$, definida pela equação $2y=x+4$. 
    \end{exercicio}
        
    \begin{exercicio}
        Dada a circunferência $c$ definida pela equação $(x-1)^2 + (y+1)^2=9$, construa as retas $r$ e $s$ tangentes
        a circunferência nos pontos $(4,-1)$ e $(1,2)$, respectivamente. Qual o ponto de intersecção das retas $r$ e $s$?
        (Dica: esboce a circunferência e as retas.)
    \end{exercicio}
        
    \begin{exercicio}
        Dada a circunferência $c$ definida pela equação $(x-1)^2 + (y+1)^2=4$, construa as retas $r$ e $s$ tangentes
        a circunferência nos pontos $(-1,-1)$ e $(1,3)$, respectivamente. Qual o ponto de intersecção das retas $r$ e $s$?
    \end{exercicio}
        
    \begin{exercicio}
        Dada a circunferência $c$ definida pela equação $(x-1)^2 + (y+1)^2=4$, construa as retas $r$ e $s$ tangentes
        a circunferência nos pontos $(-1,-1)$ e $(1,-3)$, respectivamente. Qual a distância do ponto de intersecção 
        das retas $r$ e $s$ até o centro da circunferência?
    \end{exercicio}
        
    \begin{exercicio}
        Sejam
        \begin{enumerate}[label={\roman*.}]
            \item a circunferência $u$ dada por \[(x-1)^2+(y-2)^2=4\]
            \item a circunferência $v$ dada por \[(x+2)^2=9-y^2\]
            \item a reta $r$ que passa pelos pontos \[(0,-3) \quad \text{ e } (3,-1)\]
            \item a reta $s$ dada pela equação \[2y+3x=5\]
        \end{enumerate}
        \begin{alphalist}
            \item Ilustre as formas geométricas definidas.
            \item Calcule a distância entre os centros das circunferências $u$ e $v$.
            \item Escreva a equação da reta $a$ que passa pelos centros das circunferências $u$ e $v$.
            \item Qual a declividade da reta $a$?
            \item Qual a equação da reta $r$?
            \item Qual a relação entre a reta $r$ e a reta $a$ (elas são paralelas, concorrentes, perpendiculares)?
            \item Escreva a equação da circunferência com centro na origem e que passe pelo centro da circunferência $u$.
            \item Escreva a equação da reta paralela a reta $r$ que passa pela origem.
        \end{alphalist}
    \end{exercicio}
    
    \begin{exercicio}
        Encontre uma equação para a reta $L$ que passa pelo ponto $(-2, 4)$ e satisfaz a condição dada.
        \begin{alphalist}
            \item $L$ é uma reta vertical.
            \item $L$ é uma reta horizontal.
            \item $L$ passa pelo ponto $\left(3, \sfrac{7}{2}\right)$.
            \item A intersecção $x$ de $L$ é $3$.
            \item $L$ é paralela à reta $5x - 2y = 6$.
            \item $L$ é perpendicular à reta $4x + 3y = 6$.
        \end{alphalist}
    \end{exercicio}
  
    \begin{exercicio}
        Sabendo que formula para calcular a distância entre uma reta, $r:ax+by+c=0$, e um ponto, $P(x_0,y_0)$, é
        \[d(r,P) = \frac{|ax_0+by_0+c|}{\sqrt{a^2+b^2}}\]
        Calcule a distância entre a reta e o ponto dados
        \begin{mathlist}
            \item r: 3x+4y+2=0 \) \tabto{33mm} P(-1,1)
            \item r: x+y-1=0   \) \tabto{33mm} P(2,1)
            \item r: y=2x+1    \) \tabto{33mm} P(0,-1)
            \item r: y=x       \) \tabto{33mm} P(0,1)
        \end{mathlist}
    \end{exercicio}

    \begin{exercicio}
        Sabemos que a reta tangente a uma circunferência é perpendicular ao seu raio.
        Usando essa informação, calcule a equação da reta tangente a circunferência 
        $x^2+y^2-25=0$ que passa pelo ponto $(3,4)$.
    \end{exercicio}
        
    \begin{exercicio}
        Determine se as afirmações dadas são verdadeiras ou falsas. 
        Se verdadeiras, explique porquê. Se falsas, justifique através de um exemplo.
        \begin{alphalist}
            \item O ponto $(-a, \, -b)$ é simétrico ao ponto $(a, \, b)$ em relação à origem.
            \item Se a distância entre os pontos $P_1 (a, \, b)$ e $P_ 2 (c, \, d)$ é $D$, 
                então a distância entre os pontos 
                $P_1 (a, \, b)$ e $P_3 (kc, \, kd)$, $(k \, \neq \, 0)$, é dada por $|k| D$.
            \item O círculo com equação $kx^2 \, + \, ky^2 \, = \, a^2$ situa-se dentro do 
                círculo com equação $x^2 \, + \, y^2 \, = \, a^2$, desde que $k \, > \, 1$.
            \item Sejam $(x_1, y_1)$ e $(x_2, y_2)$ dois pontos no plano $xy$. 
                Mostre que a distância entre os dois pontos é dada por
                \[d = \sqrt{(x_2 - x_1)^2 + (y_2 - y_1)^2}\]
        \end{alphalist}        
    \end{exercicio}
        
    \begin{exercicio}
        Determine o perímetro do triângulo $ABC$ que é descrito as seguintes condições
        \begin{alphalist}
            \item O vértice $A$ pertence ao eixo-$x$
            \item O vértice $B$ pertence ao eixo-$y$
            \item A reta que contêm $BC$ tem equação $x-y=0$
            \item A reta que contêm $AC$ tem equação $x+2y-3=0$
        \end{alphalist}
        \begin{answers}
            \tabto{2em} $3+\sqrt{2}+\sqrt{5}$
        \end{answers}
    \end{exercicio}
        
    \begin{exercicio}
        Determine a posição relativa de cada par de retas
        \begin{alphalist}
            \item $2x- y+3=0$  \tabto{3cm}  $4x-2y=-6$
            \item $2x- y+5=0$  \tabto{3cm}  $3x-6y=-3$
            \item $ x-2y+3=0$  \tabto{3cm}  $2x+4y+ 3=0$
            \item $2x- y+3=0$  \tabto{3cm}  $4x-2y+10=0$
            \item $2x+4y+3=0$  \tabto{3cm}  $4x-2y=-6$
        \end{alphalist}
    \end{exercicio}
        
    \begin{exercicio}
        Para quais valores de $k$ as retas $(k-1)x+6y+1=0$ e $4x+(k+1)y-1=0$ são paralelas?
    \end{exercicio}
    
    \begin{exercicio}
        Ache a equação da reta que passa pelo centro da circunferência $(x-3)^2+(y-2)^2=8$ e é perpendicular à reta $x-y-16=0$.
    \end{exercicio}
    
    \begin{exercicio}
        Um quadrado tem vértices consecutivos $A(5,0)$ e $B(-1,0)$. Determine as equações das circunferências inscrita e circunscrita ao quadrado.
    \end{exercicio}
    
    \begin{exercicio}
        Encontre a região do plano cujos pontos $P(x,y)$ satisfazem as relações $x+y\leq 2$ e $x^2+y^2\leq 16$, esboce essa região.
    \end{exercicio}
    
    \begin{exercicio}
        Resolva os sistemas de inequações
        \begin{mathlist}[1]
            \item \systeme{ x^2 + y^2 \leq 25, x^2 + y^2 \geq  4 }
            \item \systeme{ x^2 + y^2 \leq  9, x\phantom{^2} + y\phantom{^2} \leq 3 }
            \item \systeme{ x^2 + y^2 \geq  1, x^2 + y^2 \leq  9 }
            \item \systeme{ x^2 + y^2 \leq  1, x\phantom{^2} + y\phantom{^2} \geq 1 }
            \item \begin{cases}
                   x^2 + y^2     &\!\!\! \leq  4 \\
                   (x-1)^2 + y^2 &\!\!\! \leq  4
                  \end{cases}
            \item \begin{cases}
                x^2 + y^2         &\!\!\! \leq 25 \\
                (x-1)^2 + (y-1)^2 &\!\!\! \geq 16
            \end{cases}
        \end{mathlist}
    \end{exercicio}
    
    \begin{exercicio}
        Sejam os conjuntos 
        \begin{align*}
            A &= \{ (x,y) \st x^2+y^2 \leq 25 \} \\
            B &= \{ (x,y) \st (x-1)^2 + (y-1)^2 \leq 16 \} 
        \end{align*}
        Determine 
        \begin{mathlist}[2]
            \item A \cup B
            \item A \cap B
            \item A - B
            \item B - A
        \end{mathlist}
    \end{exercicio}
        
    \begin{exercicio}
        Encontre os pontos onde a circunferência $(x-1)^2+(y-2)^2=29$ encontra o eixo-$x$.
    \end{exercicio}
    
    \begin{exercicio}
        Encontre os pontos onde a circunferência $(x-1)^2+(y-2)^2=26$ encontra o eixo-$y$.
    \end{exercicio}
    
\end{exercicios}

%------------------------------------------------------------------------------%
\begin{exercicios}{Desafio}
    
    \begin{exercicio}
        Se $r$ e $s$ são retas paralelas e $AX=B$ é o sistema linear construído para 
        calcular o ponto de interseção entre elas, determine $\det A$.
    \end{exercicio}
    
\end{exercicios}
%------------------------------------------------------------------------------%

